\chapter{Sets, relations and functions}\label{ch:sets}
\index{sets}

Every object in this book is built from three primitives: the \emph{set}
(a bag of things), the \emph{relation} (a pairing rule), and the
\emph{function} (a relation that assigns each input exactly one output).
After this chapter you will be able to prove statements about all three ---
in particular the two-direction discipline behind injective/surjective ---
and you will have seen the two proof machines that power all of cardinality:
the \emph{bijection argument} (count by matching) and \emph{diagonalisation}
(show a list can never be complete). These run constantly in the counting,
probability and graph chapters.

\section{Sets, subsets and operations}\label{sec:sets}
\index{set!operations}

\begin{definition}\index{set!subset}
A \emph{set} is an unordered collection of distinct objects; $x \in A$ means
``$x$ is an element of $A$''. $A \subseteq B$ means every element of $A$ is an
element of $B$; $A = B$ means $A \subseteq B$ and $B \subseteq A$ --- that
two-way containment is the standard proof template for set equalities.
\end{definition}

The operations: union $A \cup B$, intersection $A \cap B$, difference
$A \setminus B$, complement $\overline{A}$ (inside a fixed universe $U$).
Figure~\ref{fig:venn} shows the geometry you should always draw first ---
with the warning that a Venn diagram is a \emph{guide}, not a proof.

\begin{figure}[ht]\centering
% Figure: Venn diagrams for the identities --- two-set and three-set.
% Source for Figure~\ref{fig:venn}. TikZ only; regenerable by compiling main.tex.
\begin{tikzpicture}[
  circ/.style={draw=rulegray,thick},
  lab/.style={font=\footnotesize\sffamily,text=slate},
]
% ---- two-set ----
\begin{scope}
\draw[circ] (0,0) circle (0.75);
\draw[circ] (0.85,0) circle (0.75);
\node[lab] at (-0.5,0.45) {$A$};
\node[lab] at (1.35,0.45) {$B$};
\node[lab] at (0.425,0) {\scriptsize $A\cap B$};
\node[lab] at (-0.42,-0.42) {\scriptsize $A\setminus B$};
\node[lab] at (1.27,-0.42) {\scriptsize $B\setminus A$};
\node[lab,align=center] at (0.425,-1.35) {two sets: every point\\sits in one of four regions};
\end{scope}
% ---- three-set ----
\begin{scope}[xshift=48mm]
\draw[circ] (0,0.15) circle (0.7);
\draw[circ] (-0.55,-0.45) circle (0.7);
\draw[circ] (0.55,-0.45) circle (0.7);
\node[lab] at (0,0.85) {$A$};
\node[lab] at (-1.0,-0.95) {$B$};
\node[lab] at (1.0,-0.95) {$C$};
\node[lab] at (0,-0.18) {\scriptsize $A\cap B\cap C$};
\node[lab,align=center] at (0,-1.6) {three sets: eight regions ---\\one per membership pattern};
\end{scope}
\end{tikzpicture}

\caption{Draw before you prove. Notice three sets already give eight regions
--- which is why ``the picture covers all cases'' claims die at $n=3$ sets.
Use the diagram to locate the counterexample candidate, then write the
element-chase proof.}
\label{fig:venn}
\end{figure}

\begin{example}
Prove the distributive law $A \cap (B \cup C) = (A \cap B) \cup (A \cap C)$.
\par\smallskip
\textbf{Strategy.} Equality of sets = two containments. Element-chase each
direction: take an arbitrary $x$ and unfold the definitions mechanically.
\par\smallskip
\begin{solution}
($\subseteq$) Let $x \in A \cap (B \cup C)$. Then $x \in A$ and
$x \in B \cup C$. If $x \in B$, then $x \in A \cap B \subseteq$ right side; if
$x \in C$, then $x \in A \cap C \subseteq$ right side.
($\supseteq$) Let $x \in (A \cap B) \cup (A \cap C)$. Either way $x \in A$;
and $x$ lies in $B$ or $C$, so $x \in B \cup C$. Hence
$x \in A \cap (B \cup C)$. Both containments hold, so the sets are equal.
\end{solution}
\end{example}

\begin{exercise}[easy]\label{exr:in-vs-sub}
True or false, and fix the false ones: (a) $\{1\} \in \{1,2,3\}$;
(b) $\{1\} \subseteq \{1,2,3\}$; (c) $1 \subseteq \{1,2,3\}$;
(d) $\varnothing \in \{\varnothing\}$; (e) $\varnothing \subseteq A$ for
every set $A$.
\end{exercise}

\section{Power set and Cartesian product}\label{sec:powprod}
\index{power set}\index{Cartesian product}

\begin{definition}
The \emph{power set} $\mathcal{P}(A)$ is the set of all subsets of $A$.
If $|A| = n$ then $|\mathcal{P}(A)| = 2^n$ --- each element independently
votes in or out.
\end{definition}

\begin{definition}
The \emph{Cartesian product} $A \times B = \{(a,b) : a \in A,\, b \in B\}$ ---
ordered pairs. If $|A| = m$ and $|B| = n$ then $|A \times B| = mn$. Products
iterate: $A^2 = A \times A$; $\N \times \N$ is the lattice of grid points of
Figure~\ref{fig:snake}.
\end{definition}

Ordered means $(1,2) \neq (2,1)$; a product with itself carries a diagonal
$\{(a,a)\}$ that shows up again in Section~\ref{sec:diag} for a very different
reason.

\begin{figure}[ht]\centering
% Figure: the power set of {1,2,3} as a subset lattice.
% Source for Figure~\ref{fig:lattice}. TikZ only; regenerable by compiling main.tex.
\begin{tikzpicture}[
  s/.style={draw,inner sep=3pt,font=\footnotesize,rounded corners=1pt},
  e/.style={-,rulegray},
  lab/.style={font=\footnotesize\sffamily,text=slate,align=center}
]
\node[s] (e0) at (0,3.0) {$\emptyset$};
\node[s] (s1) at (-2,1.5) {$\{1\}$};
\node[s] (s2) at (0,1.5) {$\{2\}$};
\node[s] (s3) at (2,1.5) {$\{3\}$};
\node[s] (t12) at (-2,0) {$\{1,2\}$};
\node[s] (t13) at (0,0) {$\{1,3\}$};
\node[s] (t23) at (2,0) {$\{2,3\}$};
\node[s] (full) at (0,-1.5) {$\{1,2,3\}$};
\draw[e] (e0)--(s1); \draw[e] (e0)--(s2); \draw[e] (e0)--(s3);
\draw[e] (s1)--(t12); \draw[e] (s1)--(t13);
\draw[e] (s2)--(t12); \draw[e] (s2)--(t23);
\draw[e] (s3)--(t13); \draw[e] (s3)--(t23);
\draw[e] (t12)--(full); \draw[e] (t13)--(full); \draw[e] (t23)--(full);
\node[lab,align=center]{edges add one element;\\
each level has $\binom{3}{k}$ sets ---\\
$\sum_k \binom{3}{k} = 2^3 = 8$};
\end{tikzpicture}

\caption{The power set of $\{1,2,3\}$ as a lattice. Notice: each upward edge
adds exactly one element, and level $k$ holds $\binom{3}{k}$ subsets ---
the $2^n$ count falls out of the structure, not just the formula.}
\label{fig:lattice}
\end{figure}

\begin{exercise}[easy]\label{exr:powset}
List $\mathcal{P}(\{1,2,3\})$ and verify the count $2^3$. How many subsets of
$[n]$ have even size? (Prove your count by pairing each subset with its
complement where possible, or hold it for Chapter~8.)
\end{exercise}

\section{Functions: definition, domain, image}\label{sec:functions}
\index{function}

\begin{definition}\index{domain}\index{codomain}\index{image}
A \emph{function} $f : A \to B$ assigns to every $a \in A$ exactly one
$b \in B$, written $f(a) = b$. $A$ is the \emph{domain}, $B$ the
\emph{codomain}, and $f(A) = \{f(a) : a \in A\} \subseteq B$ the
\emph{image} (or \emph{range}). Two functions are equal iff they have the
same domain, same codomain, and agree on every input.
\end{definition}

Discipline that prevents most beginner errors: codomain is a \emph{claim},
image is a \emph{fact}. ``$f : \Z \to \Z$, $f(n) = n^2$'' is a perfectly good
function whose image $\{0,1,4,9,\dots\}$ misses most of its codomain. When a
problem asks you to prove something \emph{onto}, it is the image that must be
all of the codomain --- Section~\ref{sec:injsurj} turns that into a template.

\section{Injective, surjective, bijective --- with templates}\label{sec:injsurj}
\index{injective}\index{surjective}\index{bijective}

\begin{definition}
$f : A \to B$ is \emph{injective} (one-to-one) if $f(a) = f(a') \Rightarrow a = a'$
--- distinct inputs never collide. It is \emph{surjective} (onto) if
$f(A) = B$ --- every $b \in B$ is hit. It is \emph{bijective} if both hold.
\end{definition}

Figure~\ref{fig:arrows} shows what each does to the arrows. The proof
templates are fixed --- memorise the first line of each:

\begin{rulebox}[The three templates]
\index{proof template!injective/surjective}
\begin{itemize}
\item \textbf{Injective:} ``Suppose $f(a) = f(a')$. Then \dots therefore
$a = a'$.'' (Unfold the definition of $f$, cancel carefully.)
\item \textbf{Surjective:} ``Let $b \in B$. Take $a = \langle$candidate$\rangle$.
Then $f(a) = \cdots = b$.'' (You must \emph{produce} a preimage --- guess it
by solving $f(x) = b$.)
\item \textbf{Bijective:} prove \emph{both}, or exhibit an inverse $g$ and
verify $g \circ f = \mathrm{id}_A$ and $f \circ g = \mathrm{id}_B$.
\end{itemize}
\end{rulebox}

\begin{figure}[ht]\centering
% Figure: injective / surjective / bijective arrow diagrams.
% Source for Figure~\ref{fig:arrows}. TikZ only; regenerable by compiling main.tex.
\begin{tikzpicture}[
  blob/.style={draw=rulegray,rounded corners=4pt,minimum width=11mm,minimum height=20mm},
  pt/.style={circle,fill=ink,inner sep=1.2pt},
  arr/.style={->,accent},
  lab/.style={font=\footnotesize\sffamily,text=slate,align=center},
]
% ---- panel 1: injective, not surjective ----
\begin{scope}
\node[blob] (A1) at (0,0) {};
\node[blob] (B1) at (1.7,0) {};
\node[lab] at (0,-1.35) {injective,\\not surjective};
\foreach \y [count=\i] in {0.6,0,-0.6} \node[pt] (a\i) at (0,\y) {};
\foreach \y [count=\i] in {0.75,0.25,-0.25,-0.75} \node[pt] (b\i) at (1.7,\y) {};
\draw[arr] (a1.east) -- (b1.west);
\draw[arr] (a2.east) -- (b2.west);
\draw[arr] (a3.east) -- (b4.west);
\node[font=\scriptsize,text=amber,anchor=west] at (1.75,-0.75) {missed};
\end{scope}
% ---- panel 2: surjective, not injective ----
\begin{scope}[xshift=40mm]
\node[blob] at (0,0) {};
\node[blob] at (1.7,0) {};
\node[lab] at (0.85,-1.35) {surjective,\\not injective};
\foreach \y [count=\i] in {0.75,0.25,-0.25,-0.75} \node[pt] (c\i) at (0,\y) {};
\foreach \y [count=\i] in {0.6,0,-0.6} \node[pt] (d\i) at (1.7,\y) {};
\draw[arr] (c1.east) -- (d1.west);
\draw[arr] (c2.east) -- (d1.west);
\draw[arr] (c3.east) -- (d2.west);
\draw[arr] (c4.east) -- (d3.west);
\node[font=\scriptsize,text=amber,anchor=west] at (1.75,0.6) {hit twice};
\end{scope}
% ---- panel 3: bijective ----
\begin{scope}[xshift=80mm]
\node[blob] at (0,0) {};
\node[blob] at (1.7,0) {};
\node[lab] at (0.85,-1.35) {bijective};
\foreach \y [count=\j] in {0.6,0,-0.6} \node[pt] (e\j) at (0,\y) {};
\foreach \y [count=\k] in {0.6,0,-0.6} \node[pt] (f\k) at (1.7,\y) {};
\draw[arr] (e1.east) -- (f1.west);
\draw[arr] (e2.east) -- (f2.west);
\draw[arr] (e3.east) -- (f3.west);
\end{scope}
\end{tikzpicture}

\caption{Read the arrows, not the words: injective forbids two arrows into
one target, surjective forbids a missed target, bijective forbids both.
The middle panel's ``hit twice'' mark and the left panel's ``missed'' mark
are exactly the failures each template must rule out.}
\label{fig:arrows}
\end{figure}

\begin{example}\label{ex:compinj}
Prove: the composition of two injections is injective. That is, if
$f : A \to B$ and $g : B \to C$ are injective, then $g \circ f$ is injective.
\par\smallskip
\textbf{Strategy.} Chase the definitions one layer at a time: write
$(g \circ f)(a) = (g \circ f)(a')$, peel the outer layer with $g$'s
injectivity, then the inner with $f$'s.
\par\smallskip
\begin{solution}
Suppose $(g \circ f)(a) = (g \circ f)(a')$, i.e.\ $g(f(a)) = g(f(a'))$. Since
$g$ is injective, $f(a) = f(a')$. Since $f$ is injective, $a = a'$. Hence
$g \circ f$ is injective.
\end{solution}
\end{example}

The same one-layer proof gives: composition of surjections is surjective
(exercise), and composition of bijections is bijective --- which is why
``built from bijections'' is itself a counting technique
(Section~\ref{sec:cardinality}).

\begin{exercise}[medium]\label{exr:inj}
Prove injective or give a counterexample: (a) $f:\Z\to\Z$, $f(n)=2n+1$;
(b) $f:\Z\to\Z$, $f(n)=n^2$; (c) $f:\R\to\R$, $f(x)=x^3-x$.
\end{exercise}

\section{Relations, equivalence relations, partitions}\label{sec:relations}
\index{relation}\index{equivalence relation}

\begin{definition}
A \emph{relation} on $A$ is a subset $R \subseteq A \times A$; write
$a \sim b$ (or $aRb$) for $(a,b) \in R$. $R$ is
\emph{reflexive} if $a \sim a$ for all $a$; \emph{symmetric} if
$a \sim b \Rightarrow b \sim a$; \emph{transitive} if
$a \sim b \sim c \Rightarrow a \sim c$. All three together make an
\emph{equivalence relation}; its classes
$[a] = \{b : b \sim a\}$ \emph{partition} $A$ --- every element lands in
exactly one class.
\end{definition}

\begin{example}[The parity classes]\label{ex:parityclass}
On $\Z$, let $a \sim b$ iff $a - b$ is even. Show $\sim$ is an equivalence
relation and describe its classes.
\par\smallskip
\textbf{Strategy.} Verify reflexive/symmetric/transitive straight on the
definition --- three tiny checks, no cleverness --- then read the classes off
the picture (Figure~\ref{fig:classes}).
\par\smallskip
\begin{solution}
\emph{Reflexive:} $a - a = 0$ is even. \emph{Symmetric:} if $a - b$ is even
then $b - a = -(a-b)$ is even. \emph{Transitive:} if $a - b$ and $b - c$ are
even then $a - c = (a-b) + (b-c)$ is a sum of evens, hence even. So $\sim$
is an equivalence relation. $a \sim 0$ iff $a$ is even and $a \sim 1$ iff $a$
is odd, so the classes are $[0] = \{\dots,-2,0,2,\dots\}$ and
$[1] = \{\dots,-1,1,3,\dots\}$: the even and odd integers.
\end{solution}
\end{example}

\begin{figure}[ht]\centering
% Figure: partition of Z into the two parity classes.
% Source for Figure~\ref{fig:classes}. TikZ only; regenerable by compiling main.tex.
\begin{tikzpicture}[
  ev/.style={circle,draw,fill=accentsoft,minimum size=5mm,font=\scriptsize},
  od/.style={circle,draw,fill=amberbg,minimum size=5mm,font=\scriptsize},
  bin/.style={draw=rulegray,rounded corners=3pt,minimum width=30mm,minimum height=8mm,font=\footnotesize\sffamily},
  lab/.style={font=\footnotesize\sffamily,text=slate},
]
% number-line dots -4..5
\foreach \x in {-4,-2,0,2,4}
  \node[ev] at ({(\x+4)*0.62},0) {\x};
\foreach \x in {-3,-1,1,3,5}
  \node[od] at ({(\x+4)*0.62},0) {\x};
\node[lab] at (2.8,0.55) {integers};
% bins
\node[bin] (be) at (1.2,-1.6) {$[0]$ = evens};
\node[bin] (bo) at (4.4,-1.6) {$[1]$ = odds};
\foreach \x in {0,1,2,3,4}
  \draw[->,accent] ({(\x*1.24)},-0.35) -- ({0.55+\x*0.28},-1.25);
\foreach \x in {0,1,2,3,4}
  \draw[->,amber] ({(\x*1.24)+0.62},-0.35) -- ({3.9+\x*0.28},-1.25);
\node[lab] at (2.8,-2.5) {every integer lands in exactly one class};
\end{tikzpicture}

\caption{An equivalence relation sorts its set into bins. Notice nothing is
left over and nothing is in two bins --- ``partition'' is not a slogan, it is
the theorem that classes never overlap and never miss anyone.}
\label{fig:classes}
\end{figure}

This is the trailer for Chapter~6: ``same parity'' is congruence mod $2$, and
mod $m$ will give $m$ classes with the same three-line verification.

\begin{exercise}[medium]\label{exr:rel}
Decide which of reflexive/symmetric/transitive each relation on $\Z$ has:
(a) $a \sim b$ iff $a \le b$; (b) $a \sim b$ iff $|a - b| \le 1$;
(c) $a \sim b$ iff $ab \ge 0$. Which are equivalence relations?
\end{exercise}

\section{Cardinality: bijections count, diagonalisation escapes}\label{sec:cardinality}\label{sec:diag}
\index{cardinality}\index{countable}

Two finite sets have the same size iff a bijection exists between them ---
that is the \emph{bijection principle}, and Chapter~8 spends a book on it.
Applied to infinite sets it gives the first surprise:

\begin{definition}
A set $A$ is \emph{countably infinite} if a bijection $\N \to A$ exists;
\emph{countable} means finite or countably infinite; otherwise
\emph{uncountable}.
\end{definition}

$\Z$ is countable (interleave $0,1,-1,2,-2,\dots$). $\Q$ is countable --- the
snake of the next example is the standard proof sketch.

\begin{example}\label{ex:nn}
Show $\N \times \N$ is countable.
\par\smallskip
\textbf{Strategy.} Draw a snake through the grid --- a picture argument made
precise. Visit anti-diagonals $x + y = k$ for $k = 2,3,4,\dots$; each diagonal
is finite, so every lattice point is reached at a finite step.
\par\smallskip
\begin{solution}
List pairs in order of increasing $x + y$, traversing each diagonal in the
direction that continues the path --- alternately decreasing-$x$ and
increasing-$x$, which is what keeps the trail continuous:
\[ (1,1),\ (2,1),\ (1,2),\ (1,3),\ (2,2),\ (3,1),\ (4,1),\ \dots \]
the path in
Figure~\ref{fig:snake}. Every $(x,y)$ lies on the finite diagonal $x+y=x_0+y_0$
and is enumerated at a finite position, so the list is a bijection
$\N \to \N \times \N$. Hence $\N \times \N$ is countable.
\end{solution}
\end{example}

\begin{figure}[ht]\centering
% Figure: the snake through N x N (N = {1,2,...}) --- one path that visits every lattice point.
% Source for Figure~\ref{fig:snake}. TikZ only; regenerable by compiling main.tex.
\begin{tikzpicture}[
  pt/.style={circle,fill=ink,inner sep=1.1pt},
  arr/.style={->,accent,thick},
  lab/.style={font=\scriptsize,text=slate},
]
\foreach \i in {0,...,4}
  \foreach \j in {0,...,4}
    \node[pt] (p\i\j) at (\j*0.72,-\i*0.72) {};
% axis labels
\node[lab,anchor=east] at (-0.15,0) {$1$};
\node[lab,anchor=east] at (-0.15,-0.72) {$2$};
\node[lab,anchor=east] at (-0.15,-1.44) {$3$};
\node[lab,anchor=north] at (0,0.3) {$1$};
\node[lab,anchor=north] at (0.72,0.3) {$2$};
\node[lab,anchor=north] at (1.44,0.3) {$3$};
\node[lab,anchor=north,font=\footnotesize,text=slate] at (1.44,0.75) {$x$};
\node[lab,anchor=east,font=\footnotesize,text=slate] at (-0.55,-1.44) {$y$};
% snake path: visit order
\draw[arr] (p00.east) -- (p01.west);
\draw[arr] (p01.south) -- (p10.north);
\draw[arr] (p10.east) -- (p20.west);
\draw[arr] (p20.north east) -- (p11.south west);
\draw[arr] (p11.north west) -- (p02.south east);
\draw[arr] (p02.east) -- (p03.west);
\draw[arr] (p03.south) -- (p12.north);
\draw[arr] (p12.south east) -- (p21.north west);
\draw[arr] (p21.south east) -- (p30.north west);
\draw[arr] (p30.east) -- (p40.west);
\draw[arr] (p40.north) -- (p31.south);
\draw[arr] (p31.north west) -- (p22.south east);
\draw[arr] (p22.north west) -- (p13.south east);
\draw[arr] (p13.north west) -- (p04.south east);
% step numbers
\node[lab,anchor=south] at (0,0.28) {\!1};
\node[lab,anchor=south] at (0.72,0.28) {\!2};
\node[lab,anchor=south west] at (0.02,-0.7) {3\ };
\node[lab,anchor=north west] at (-0.02,-1.4) {4\ };
\node[lab,anchor=west] at (1.5,-0.74) {5};
\node[lab,anchor=south] at (1.44,0.26) {6};
\end{tikzpicture}

\caption{The snake enumeration of $\N \times \N$ (here $\N =
\{1,2,3,\dots\}$). Notice why it works: every
anti-diagonal $x+y=k$ is \emph{finite}, so the path reaches any given point
after finitely many steps --- ``eventually covers everything'' is the whole
content of countable.}
\label{fig:snake}
\end{figure}

And the second surprise --- there are strictly more reals:

\begin{example}[Cantor diagonalisation]\label{ex:cantor}
\index{diagonalisation}
Prove: $[0,1)$ is uncountable --- no list of reals in $[0,1)$ is complete.
\par\smallskip
\textbf{Strategy.} Given any candidate list, \emph{manufacture the missing
witness}: build a number whose $n$-th decimal digit differs from the $n$-th
digit of the $n$-th row. It differs from every row, so it was never listed.
\par\smallskip
\begin{solution}
Suppose $r_1, r_2, r_3, \dots$ is any list of numbers in $[0,1)$, each written
as decimals $0.d_{i1}d_{i2}d_{i3}\dots$ (forbid the all-9s tail so every
number has one expansion). Build $x = 0.e_1e_2e_3\dots$ by the rule
$e_i = 4$ if $d_{ii} = 5$, else $e_i = 5$. For every $i$, $x$ differs from
$r_i$ in position $i$, so $x \neq r_i$; and $x \in [0,1)$ since its digits
are 4s and 5s --- never the ambiguous 9-tail. So $x$ is a number in $[0,1)$
absent from the list. No list is complete; $[0,1)$ is uncountable.
\end{solution}
\end{example}

\begin{figure}[ht]\centering
% Figure: Cantor diagonalisation --- flip each marked digit, escape the list.
% Source for Figure~\ref{fig:diagonal}. TikZ only; regenerable by compiling main.tex.
\begin{tikzpicture}[
  dg/.style={font=\footnotesize\ttfamily,text=ink},
  hit/.style={draw=amber,thick,rounded corners=1pt},
  lab/.style={font=\footnotesize\sffamily,text=slate},
  new/.style={font=\footnotesize\ttfamily,text=accent},
]
\def\cw{0.21}
% the table (digits monospaced); rows 0..7, decimal digits at columns 2..11
\foreach \r/\row in {
  0/{0.5000000000},
  1/{0.3333333333},
  2/{0.1415926535},
  3/{0.7182818284},
  4/{0.4142135623},
  5/{0.6180339887},
  6/{0.1010010001},
  7/{0.2500000000}}
  \node[dg,anchor=west] at (0,-\r*0.34) {\row};
% amber boxes on the diagonal digit of each row: column index = 2 + r
\foreach \r in {0,...,7}
  \node[hit,minimum width=2.6mm,minimum height=3mm] at ({(2+\r)*\cw + 0.018},-\r*0.34-0.02) {};
\node[lab,anchor=west] at (2.7,-0.1) {each row's marked digit};
\draw[->,slate] (2.75,-0.45) -- (1.28,-1.35);
% constructed row
\node[new,anchor=west] at (0,-2.95) {0.45555555...};
\node[lab,anchor=west] at (2.2,-2.98) {flipped each marked digit --- absent from every row};
\end{tikzpicture}

\caption{The flip rule at work. Notice two engineering details that make the
argument watertight: the constructed digits are only 4s and 5s (no 9-tail, so
$x$'s expansion is unique), and each row is attacked at \emph{its own}
position --- the new number escapes all rows simultaneously.}
\label{fig:diagonal}
\end{figure}

\begin{checkpoint}[Predict before you run]
Listing~\ref{lst:diag} runs the construction on eight decimals. Before
running: predict the constructed number's first two digits, and predict which
row it could \emph{still} equal even though it differs at the marked
position. (Trick: none --- differing in one decimal position is enough when
expansions are unique.)
\end{checkpoint}

\begin{lstlisting}[language=Python,caption={\texttt{diagonal\_demo.py}: the
manufacture-the-witness step, mechanised (full script in
\texttt{code/diagonal\_demo.py}).},label={lst:diag}]
def flip(digit):                 # never emits 9 -> unique expansion
    return "4" if digit == "5" else "5"

for i, row in enumerate(ROWS):
    ...                          # print row with position i+2 marked
    digits.append(flip(row[i+2]))
built = "0." + "".join(digits) + "..."
\end{lstlisting}
Real output:
\begin{lstlisting}[numbers=none]
0.[5]000000000
0.3[3]33333333
0.14[1]5926535
0.718[2]818284
0.4142[1]35623
0.61803[3]9887
0.101001[0]001
0.2500000[0]00
constructed: 0.45555555...
differs from every row at its marked position
\end{lstlisting}
Eight rows, eight marked positions, one number that is provably on none of
them. Infinite lists differ only in scale, not in kind.

\begin{exercise}[easy]\label{exr:countable}
Show each set is countable by describing the enumeration:
(a) $\Z$; (b) the even integers; (c) $\N \times \{0,1\}$.
\end{exercise}

\section{The set $[n]$ and indexing conventions}\label{sec:bracketn}
\index{[n]@$[n]$}

\begin{definition}
$[n] = \{1, 2, \dots, n\}$, the standard $n$-element set; $[0] = \varnothing$.
Any set of size $n$ is ``a copy of $[n]$'' via some bijection --- that is the
content of counting.
\end{definition}

Conventions this book keeps: indices run $1..n$ unless declared otherwise;
$a_1, \dots, a_n$ is a \emph{sequence} (ordered, may repeat) while
$\{a_1, \dots, a_n\}$ is a set (unordered, distinct); $x \in [n]$ means
$1 \le x \le n$ integer. $\N = \{1, 2, 3, \dots\}$ throughout --- the olympiad
convention of the notation guide --- with $\Z_{\ge 0}$ written explicitly
when $0$ must be included.

\section{Chapter summary}
\begin{table}[ht]\centering
\begin{tabularx}{\linewidth}{@{}lYY@{}}
\toprule
\tablehead\thead{To prove\dots} & \thead{Template} & \thead{First line}\\
\midrule
$A = B$ & two containments & ``Let $x \in A$'' / ``Let $x \in B$''\\
$f$ injective & unfold collision & ``Suppose $f(a)=f(a')$''\\
$f$ surjective & produce preimage & ``Let $b \in B$; take $a = \dots$''\\
$f$ bijective & both, or inverse & verify $g\circ f$, $f\circ g$ identities\\
$\sim$ equivalence & R/S/T on definition & $a\sim a$; flip; add\\
$A$ countable & enumerate & list that reaches every element\\
$A$ uncountable & diagonalise & ``Suppose $r_1, r_2, \dots$ is a list''\\
count finite sets & bijection & match the objects, count $[n]$\\
\bottomrule
\end{tabularx}
\caption{Every claim shape in this chapter has a named template and a
required first line. Write the first line before you think --- it is the
part graders grade.}
\label{tab:sets}
\end{table}

\section{Common mistakes}
\index{common mistakes}
\begin{itemize}
\item \textbf{Confusing $\in$ with $\subseteq$.} $1 \in \{1,2\}$ but
$\{1\} \subseteq \{1,2\}$; the wrapped version is a different object.
Check every statement's type before arguing about it.
\item \textbf{Proving a bijection by one direction.} Injective-only and
surjective-only examples are in Figure~\ref{fig:arrows}; one property says
nothing about the other. Prove both, or exhibit the inverse.
\item \textbf{``Uncountable'' $\neq$ ``infinite''.} $\Z$ and $\Q$ are infinite
\emph{and} countable; the failure mode is assuming a list exists because the
set is ``just'' infinite. Diagonalisation is precisely the argument that no
list suffices for $[0,1)$.
\end{itemize}

\section{Additional exercises}
\begin{exercise}[easy]
Prove by element-chase: $A \cup (B \cap C) = (A \cup B) \cap (A \cup C)$.
\end{exercise}
\begin{exercise}[easy]
Prove De Morgan: $\overline{A \cup B} = \overline{A} \cap \overline{B}$.
(Draw Figure~\ref{fig:venn}'s left panel first.)
\end{exercise}
\begin{exercise}[easy]
Give an example of a function $f : \Z \to \Z$ that is injective but not
surjective, and one that is surjective but not injective --- and prove each
property you claim.
\end{exercise}
\begin{exercise}[easy]
On $\Z$, is $a \sim b$ iff $a \mid b$ an equivalence relation? Which of
reflexive/symmetric/transitive hold?
\end{exercise}
\begin{exercise}[medium]
Prove: the composition of two surjections is surjective. (Mirror
Example~\ref{ex:compinj}; the harder half is producing the preimage.)
\end{exercise}
\begin{exercise}[medium]
Show the positive rationals $\Q_{>0}$ are countable. (Snake the pairs
$(p,q)$, skip non-reduced fractions --- explain why skipping is allowed.)
\end{exercise}
\begin{exercise}[medium]
Let $a \sim b$ on $\Z$ iff $a - b$ is divisible by $3$. Prove it is an
equivalence relation and describe all its classes.
\end{exercise}
\begin{exercise}[medium]
Prove: if $f : A \to B$ and $g : B \to C$ with $g \circ f$ injective, then
$f$ is injective. Does $g$ have to be injective? Prove or counterexample.
\end{exercise}
\begin{exercise}[hard]
Adapt diagonalisation to subsets: prove $\mathcal{P}(\N)$ is uncountable by
supposing a list $S_1, S_2, \dots$ and building a subset that differs from
every $S_i$ on membership of $i$.
\end{exercise}
\begin{exercise}[hard]
Prove Cantor's theorem: for \emph{every} set $A$, no function
$f : A \to \mathcal{P}(A)$ is surjective. (Consider the diagonal set
$D = \{a \in A : a \notin f(a)\}$ --- Russell's paradox in one line.)
\end{exercise}
\begin{exercise}[hard]
Extend \texttt{code/diagonal\_demo.py} so it works for base-3 expansions
(digits 0--2), choosing a flip rule that still avoids ambiguity. Then prove
from your construction that the numbers in $[0,1)$ with \emph{only} digits
$\{0,2\}$ in base 3 (the Cantor set) are uncountable.
\end{exercise}
\begin{exercise}[easy]
Let $A = \{1,2,3,4\}$, $B = \{2,4,6\}$. Compute $A\cup B$, $A\cap B$,
$A\setminus B$, $A\times B$'s cardinality, and $|\mathcal{P}(A\cap B)|$.
\end{exercise}
\begin{exercise}[medium]
Prove or disprove: $f(n) = 2n+1$ as a function $\Z\to\Z$ is (i) injective,
(ii) surjective, (iii) bijective. If any fails, exhibit the counterexample.
Then restrict the codomain to the odds and redo all three.
\end{exercise}
\begin{exercise}[medium]
On $\Z$, define $a \sim b$ iff $a - b$ is divisible by $5$. Prove this is an
equivalence relation, list its classes, and say in one sentence why a
partition and an equivalence relation are the same object.
\end{exercise}
\begin{exercise}[hard]
Prove that the interval $(0,1)$ and the open ray $(0,\infty)$ have the same
cardinality by writing an explicit bijection (one of $x \mapsto x/(1-x)$ or
its cousin) and checking it is injective and surjective --- not by waving at
``infinite means infinite''.
\end{exercise}

\section*{Further reading}
\begin{itemize}
\item R.~Hammack, \emph{Book of Proof}, 3rd ed., Ch.~1, 11--13 ---
\url{https://richardhammack.github.io/BookOfProof/} --- free PDF; sets,
functions and cardinality with the two-containment and injective/surjective
templates this chapter uses.
\item E.~Lehman, F.~T.~Leighton, A.~R.~Meyer, \emph{Mathematics for Computer
Science}, Ch.~4, 8 ---
\url{https://courses.csail.mit.edu/6.042/spring18/mcs.pdf} --- free PDF;
the countability order and light Cantor treatment of \S\ref{sec:cardinality}.
\item E.~Chen, \emph{Introduction to Functional Equations} ---
injective, surjective and bijective restated in solve-for-the-witness
form; the bridge from this chapter to olympiad functional equations:

\url{https://web.evanchen.cc/handouts/FuncEq-Intro/FuncEq-Intro.pdf}
\item Stanford CS103, \emph{Mathematical Foundations of Computing} ---
\url{https://web.stanford.edu/class/cs103/} --- free reader; its
set-theory and functions units drill element chasing to fluency.
\item The Open Logic Project ---
\url{https://openlogicproject.org/} --- free texts; for the
axiomatic-set-theory view beneath the naive one taught here.
\item MIT OpenCourseWare, \emph{6.042J} (Fall 2010) ---
\url{https://ocw.mit.edu/6-042JF10} --- the sets and cardinality
lectures behind \S\ref{sec:cardinality}.
\end{itemize}

\begin{rulebox}[Watch instead]
\begin{itemize}
\item \textbf{TrevTutor --- ``INTRODUCTION to SET THEORY''}, 16:38,
\url{https://www.youtube.com/watch?v=tyDKR4FG3Yw} --- covers
\S\ref{sec:sets}--\ref{sec:powprod}: sets, membership, subsets,
power sets. (Verified 2026-10-11.)
\item \textbf{TrevTutor --- ``INJECTIVE, SURJECTIVE, and BIJECTIVE
FUNCTIONS''}, 17:14,
\url{https://www.youtube.com/watch?v=bZred_Ksz2k} --- covers
\S\ref{sec:injsurj}: the three properties with arrow-diagram examples
matching Figure~\ref{fig:arrows}. (Verified 2026-10-11.)
\end{itemize}
\end{rulebox}

\chapternotes
The element-chase template, the two-containment proof of set equality, and
the treatment of relations/partitions follow Velleman's chapters 4--5 and
8~\cite{velleman}; the injective/surjective templates and the composition
exercises are Hammack's chapter 11--12 approach~\cite{hammack}, with the
definition conventions of Chen's \emph{Introduction to Functional
Equations}~\cite{chen-funceq} (injections, surjections and bijections
stated in solve-for-the-witness form). MIT 6.042J's chapters on sets and cardinality
supplied the countability order of Section~\ref{sec:cardinality} and the
light Cantor treatment~\cite{mcs}; diagonalisation as ``manufacture the
missing witness'' is the standard presentation~\cite{hammack,mcs}. The
counting-flavoured contest problems this chapter feeds are
AIME 2024 II P6, AIME 2023 II P11, AIME 2022 I P12, AIME 2010 I P7 and
AMC 12 2023 A P24, all on the AoPS archive~\cite{aops-amc}.
